\subsection{Phân hệ quản lý tài liệu và biểu mẫu}

Phân hệ quản lý tài liệu và biểu mẫu được truy cập thông qua route \texttt{/docs}. Màn hình này hiển thị danh sách tài liệu biểu mẫu với các cột như Document Name, Description, Tags, File Type, File Size, Created At và Actions. Các tài liệu quan sát được gồm nhiều loại biểu mẫu như đơn hủy học phần, giấy giới thiệu thực tập, đơn xin đăng ký đồ án lần 2, đơn xin cấp lại thẻ bảo hiểm y tế, đơn xin bảo lưu điểm thi lại, đơn xin phúc khảo bài thi, đơn xin hoãn thi, thanh toán ra trường và đơn xin thôi học.

\subsubsection{Mục đích của quản lý tài liệu và biểu mẫu}

Chức năng này giúp nhà trường tập trung quản lý các biểu mẫu thường dùng để người dùng có thể tra cứu, tải xuống hoặc được chatbot hướng dẫn sử dụng khi cần. Trong bối cảnh tuyển sinh và đào tạo, biểu mẫu là nhóm tài liệu có tần suất sử dụng cao, đặc biệt đối với sinh viên, thí sinh và cán bộ xử lý hồ sơ.

Việc đưa biểu mẫu vào hệ thống quản trị giúp giảm thao tác thủ công, hạn chế tình trạng sử dụng biểu mẫu cũ và tạo nền tảng để chatbot có thể trả lời các câu hỏi như: ``Em muốn xin hoãn thi thì dùng mẫu nào?'', ``Mẫu đơn xin phúc khảo ở đâu?'' hoặc ``Có giấy giới thiệu thực tập không?''.

\subsubsection{Thông tin tài liệu cần quản lý}

Mỗi tài liệu hoặc biểu mẫu cần lưu các thông tin chính sau:

\begin{itemize}
    \item Tên tài liệu.
    \item Mô tả tài liệu.
    \item Danh sách tag.
    \item Loại file.
    \item Dung lượng file.
    \item Thời điểm tạo.
    \item Trạng thái public/private nếu có.
    \item Các thao tác quản trị như xem, sửa, tải xuống, khôi phục hoặc xóa mềm.
\end{itemize}

Các trường này phù hợp với các cột đã quan sát trên màn hình \texttt{/docs}. Trong đó, tag giúp phân loại tài liệu theo nhóm nghiệp vụ, loại file giúp người dùng biết định dạng tài liệu và dung lượng file giúp kiểm soát lưu trữ.

\subsubsection{Tìm kiếm, lọc và thao tác tài liệu}

Màn hình quản lý tài liệu cần hỗ trợ tìm kiếm theo tên tài liệu, mô tả và tag. Ngoài ra, hệ thống cần cho phép lọc theo loại file, trạng thái public/private hoặc thời gian tạo. Đối với danh sách tài liệu lớn, cần có phân trang để đảm bảo giao diện dễ sử dụng.

Các thao tác cơ bản gồm:

\begin{itemize}
    \item Thêm tài liệu mới.
    \item Cập nhật thông tin mô tả hoặc tag.
    \item Tải xuống tài liệu.
    \item Xóa mềm tài liệu.
    \item Khôi phục tài liệu đã xóa mềm.
\end{itemize}

Việc xuất hiện quyền \texttt{POST /api/v1/documents/:id/restore} và \texttt{DELETE /api/v1/documents/:id/soft} cho thấy hệ thống đã tính đến cơ chế khôi phục và xóa mềm tài liệu. Đây là thiết kế phù hợp vì tài liệu quản trị không nên bị xóa vĩnh viễn ngay lập tức, nhằm đảm bảo khả năng truy vết và phục hồi khi thao tác nhầm.
